\section{Query-Specific Window Learning}
\subsection{Two true relations, two different answers}
In one recording $x$, RelTwin arranges the same ESC-50 excerpts~\cite{piczak2015esc}, $A$ and $B$, as $A\!\rightarrow\!B$ and $B\!\rightarrow\!A$. Their local targets $W_+$ and $W_-$ run from the first excerpt's onset to the second's offset, including the gap. Queries $q_+$ and $q_-$ request the respective orders. Reusing the excerpts holds acoustic content fixed across windows, so their labels depend on the queried order: $W_-$ is positive for $q_-$ and negative for $q_+$.

We alternate which target order appears first in the recording. Category rehearsal targets the union of both occurrences of $A$ or $B$, retaining category-grounding supervision alongside relation selection.

\subsection{Contrast timestamp answers during training}
RelTwin scores complete timestamp answers under each query and trains their relative preference. We serialize an interval set $S$ in the backbone's native format: comma-separated spans for SpotSound-A and temporal tokens for TimeAudio. Answer tokens $y(S)$ exclude the prompt and receive the length-normalized score
\begin{equation}
s_\theta(x,q,S)=\frac{1}{|y(S)|}\sum_t\log P_\theta(y_t\mid x,q,y_{<t}).
\label{eq:score}
\end{equation}
The $2\times2$ matrix $M_{ij}=s_\theta(x,q_i,W_j)$ for $i,j\in\{+,-\}$ indexes queries by rows and answers by columns; diagonal entries correspond to correct pairs (Fig.~\ref{fig:pair}b).
Sequence supervision minimizes $\mathcal L_{\rm seq}=-(M_{++}+M_{--})/2$. Candidate supervision explicitly trains the preferences $M_{++}>M_{+-}$ and $M_{--}>M_{-+}$. With row-wise probabilities $p_i(j)=\operatorname{softmax}_j(M_{ij}/\tau)$, $\tau>0$, its loss is
\begin{equation}
\mathcal L_{\rm cand}=-\tfrac12[\log p_+(+)+\log p_-(-)].
\label{eq:candidate}
\end{equation}
Query reversal exchanges the target assignment while preserving the recording. With label permutation $\Pi$, query-exchange consistency $\mathcal L_{\rm JS}=\allowbreak\operatorname{JS}(p_+,\Pi p_-)$ couples the paired preferences after relabeling. For rehearsal example $(x_r,q_r,S_r)$, define $\mathcal R_1=-s_\theta(x_r,q_r,S_r)$. Stage~1 optimizes
\begin{equation}
\mathcal L_1=\mathcal L_{\rm seq}+\lambda(\mathcal L_{\rm cand}
+\gamma\mathcal L_{\rm JS})+\beta\mathcal R_1.
\label{eq:loss}
\end{equation}

\textbf{Set-level preference learning.}
Selecting the relevant window also requires controlling its boundaries, coverage, and event count. Stage~2 extends answer competition to six interval sets: $W_\pm$, their expanded versions, their enclosing span, and their union. Quality axes are set-IoU, precision--recall margins above IoU 0.3/0.5, count agreement, and duration-normalized boundary accuracy. Utility $v_{ij}$ averages mean quality and net Pareto dominance (number dominated minus dominating, normalized by candidate count minus one). This yields soft targets $t_i=\operatorname{softmax}(v_i/\tau_q)$ and $\mathcal L_{\rm set}=-\tfrac12\sum_{i,j}t_i(j)\log p_i(j)$ over six answers, replacing $\mathcal L_{\rm cand}$ in Eq.~\ref{eq:loss}. Exchange $\Pi$ swaps targets and their expansions, fixing the span and union.

Replay candidates perturb boundaries, add/drop events, merge spans, split the longest interval, or cover the recording. Stage~2 uses $\mathcal R_2=\mathcal R_1+\lambda\mathcal L_{\rm set,r}+\kappa\operatorname{KL}(p_0\Vert p_r)$, with frozen base-model distribution $p_0$; $\mathcal L_{\rm set,r}$ is the same listwise loss for one replay query. Inference directly generates timestamps.
